\section{The Four Pillars of Object Orientation}
\label{sec:lld-four-pillars}

\problemstatement{Explain encapsulation, abstraction, inheritance, and
polymorphism -- what each means, why it exists, and how it looks in C++ --
since every later pattern is built from them.}

\begin{patternbox}[What Each Pillar Is Really For]
\begin{itemize}
  \item \textbf{Encapsulation} -- bundle data with the methods that guard
        it, and hide the data behind those methods, so invariants cannot
        be broken from outside.
  \item \textbf{Abstraction} -- expose \emph{what} an object does, hide
        \emph{how}, so callers depend on a stable interface, not
        internals.
  \item \textbf{Inheritance} -- reuse and specialise a base type,
        modelling an ``is-a'' relationship.
  \item \textbf{Polymorphism} -- let one interface stand for many
        concrete behaviours chosen at run time, so calling code needn't
        know the exact subtype.
\end{itemize}
\end{patternbox}

\subsection*{Encapsulation and abstraction: two sides of the same wall}

Encapsulation is about \emph{protection}: making fields \texttt{private}
and offering controlled access so no caller can leave an object in an
invalid state (a bank balance that only \texttt{deposit}/\texttt{withdraw}
can change, never a raw assignment). Abstraction is about
\emph{simplification}: the caller sees \texttt{account.withdraw(100)} and
neither knows nor cares how overdraft rules are enforced. Encapsulation
hides data; abstraction hides implementation. In practice the same
\texttt{private} fields plus a small public interface deliver both.

\begin{ideabox}[Program to an Interface, Not an Implementation]
The single most valuable habit the four pillars enable: callers should
depend on an \emph{abstract type} (a base class or interface), so the
concrete class behind it can be swapped without touching the caller. This
one idea is the seed of polymorphism, the Strategy pattern, dependency
injection, and testability.
\end{ideabox}

\subsection*{Inheritance and polymorphism: one call, many behaviours}

Inheritance lets a \texttt{Dog} reuse and extend an \texttt{Animal}.
Polymorphism lets a function that accepts an \texttt{Animal*} call
\texttt{speak()} and get barking or meowing depending on the real object,
resolved at run time through a \texttt{virtual} function. This run-time
dispatch is what allows open-ended designs: new animal types slot in
without changing the function that iterates over animals.

\subsection*{C++ illustration}

\begin{lstlisting}
#include <bits/stdc++.h>
using namespace std;

// Encapsulation + abstraction: data hidden, invariant enforced.
class Account {
    double balance = 0;                 // private: cannot be corrupted
public:
    void deposit(double amt) { if (amt > 0) balance += amt; }
    bool withdraw(double amt) {
        if (amt > 0 && amt <= balance) { balance -= amt; return true; }
        return false;                   // overdraft rule hidden inside
    }
    double getBalance() const { return balance; }
};

// Inheritance + polymorphism: one interface, many behaviours.
class Animal {
public:
    virtual string speak() const = 0;   // pure virtual: abstraction
    virtual ~Animal() = default;
};
class Dog : public Animal {
public:
    string speak() const override { return "Woof"; }
};
class Cat : public Animal {
public:
    string speak() const override { return "Meow"; }
};

int main() {
    Account a;
    a.deposit(100);
    cout << "Withdrew 40? " << a.withdraw(40) << ", balance "
         << a.getBalance() << "\n";

    vector<unique_ptr<Animal>> zoo;
    zoo.push_back(make_unique<Dog>());
    zoo.push_back(make_unique<Cat>());
    for (auto& animal : zoo)            // caller knows only "Animal"
        cout << animal->speak() << " ";
    cout << "\n";
    return 0;
}
\end{lstlisting}

\begin{takeawaybox}
Encapsulation guards invariants, abstraction hides implementation,
inheritance reuses, and polymorphism defers the choice of behaviour to run
time. The recurring pay-off is the same: \emph{callers depend on a stable
abstraction, so concrete details can change without a ripple}. Every
pattern in the next chapter is a specific, named way of exploiting that.
\end{takeawaybox}
